\section{Discussão}
\label{sec:discussao}

Os resultados obtidos mostram que a taxa estimada de desempenho aumentou conforme a ordem das matrizes cresceu, como sintetizado na Figura~\ref{fig:benchmark-gflops}. A execução usou 16 threads no OpenBLAS. A média passou de 28,595 GFLOPS, na ordem 512, para 94,367 GFLOPS, na ordem 2048. O ganho, porém, não foi uniforme: a passagem de 512 para 1024 elevou a média apenas 12,2\%, enquanto a passagem de 1024 para 2048 apresentou ganho aproximado de 194\%.

Esse comportamento sugere que matrizes pequenas e médias ainda sofrem influência de custos fixos de chamada, alocação, agendamento de threads e preparação dos dados. Na ordem 2048, a carga passa a ter trabalho computacional suficiente para aproveitar melhor a rotina de multiplicação de matrizes do NumPy/OpenBLAS. A relação com a hierarquia de cache também ajuda a interpretar o salto: o processador reportou L2 total de 10 MiB e L3 de 18 MiB, e matrizes maiores aumentam a pressão sobre cache e memória, tornando mais visível o efeito de bloqueio, movimentação de dados e paralelismo explorado pela biblioteca.

A variação entre repetições também mudou conforme o tamanho da matriz. Na ordem 1024, o desvio-padrão foi de 15,385 GFLOPS para média de 32,082 GFLOPS, coeficiente de variação de aproximadamente 48\%. Essa dispersão é alta e deve ser interpretada como ruído experimental relevante. Uma causa provável é a execução em WSL2, sujeita a interferências do sistema hospedeiro, agendamento de threads, uso de swap e variações térmicas. O teste $t$ de Welch não indicou diferença significativa entre 512 e 1024, mas indicou diferença entre 1024 e 2048, considerando $\alpha=0{,}05$.

Embora o experimento não execute um modelo generativo completo, a multiplicação de matrizes é uma operação central em redes neurais modernas. Assim, o benchmark permite observar, em escala controlada, a relação entre carga computacional, biblioteca numérica e infraestrutura disponível. A máquina avaliada possui CPU Intel Core i7-1360P, 32 GB de RAM e GPU integrada Intel Iris Xe, mas a execução foi realizada em CPU, no Ubuntu 24.04 LTS via WSL2, sem aceleração explícita por GPU.

A comparação com o NVIDIA DGX H100 evidencia a diferença de escala entre uma máquina local e uma infraestrutura especializada para IA. Enquanto o ambiente local oferece uma linha de base acessível para experimentação, sistemas dedicados combinam múltiplas GPUs, grande volume de memória de GPU, rede de alta velocidade e pilha de software otimizada.

As principais limitações são a execução em apenas uma máquina, o uso de benchmark sintético e a ausência de medição direta de energia, temperatura e GPU. O ambiente Linux também não foi completamente isolado: a execução ocorreu em WSL2, escolha feita pela acessibilidade educacional, com processos do sistema hospedeiro, uso de swap e variações térmicas não controlados de forma rígida. Ainda assim, os dados cumprem o objetivo de demonstrar uma metodologia inicial de avaliação técnica e fornecem base para experimentos futuros com modelos generativos reais, aceleração por GPU ou comparação entre ambientes.
